\section{Performance Benchmarking}\label{sec:evaluation-benchmarking}

A defining characteristic of the proposed DSL compiler is its "zero-runtime" architecture, which shifts the computational burden entirely to the compilation phase. To evaluate the viability of this approach, it is essential to analyze the compiler's execution speed, particularly as the complexity of the musical composition scales.

\subsection{Theoretical Complexity}

The compilation pipeline is composed of several sequential passes, each with a distinct computational complexity:
\begin{enumerate}
    \item \textbf{Frontend and Semantic Analysis:} The complexity of parsing and type-checking is proportional to the number of lines of source code, $O(S)$, where $S$ is the size of the abstract syntax tree. Because these passes do not unroll loops or instantiate patterns multiple times, they remain extremely fast regardless of the composition's length.
    \item \textbf{Lowering Engine:} The unrolling process has a time complexity of $O(E)$, where $E$ is the total number of generated \texttt{NoteEvent}s. A simple \texttt{loop (1000)} statement will execute its body 1,000 times during this phase, meaning the lowering time is proportional to the density of the final timeline, not the brevity of the source code.
    \item \textbf{MIDI Backend:} The backend must sort the absolute events by their tick values before computing delta-times. This stable sort introduces a time complexity of $O(E \log E)$. Serializing the events into the binary file requires a final linear pass, $O(E)$.
\end{enumerate}
Therefore, the overall time complexity of the compilation process is bounded by the sorting phase in the backend: $O(E \log E)$.

\subsection{Execution Speed and Real-Time Viability}

Because the compiler is implemented in highly optimized C++23 without reliance on garbage collection or dynamic dispatch for AST traversal, the constant factors in the time complexity are exceptionally small.
In practice, the compiler generates compositions containing tens of thousands of note events in a fraction of a second.

This high-performance metric is critical for the utility of the DSL in a real-world scenario. The interactive playground environment (discussed in previous chapters) relies on the compiler's ability to re-compile the entire source file on every keystroke. Because the execution time for typical musical pieces (ranging from hundreds to thousands of events) sits well below the threshold of human perception (typically $<100$ milliseconds), the composer experiences real-time auditory and visual feedback without the need for an interpreted runtime.

\subsection{Memory Footprint}

The value-oriented design of the AST and IR ensures that the compiler's memory footprint remains minimal. The use of \texttt{std::variant} guarantees spatial locality in the CPU cache during AST traversal. The intermediate representation holds only the necessary atomic data (pitch, start beat, duration, velocity) packed tightly into \texttt{struct}s. Consequently, even a composition unrolled into millions of events will consume only a few megabytes of RAM, allowing the compiler to run seamlessly on low-end hardware or within memory-constrained Docker containers.
